% =============================================================================
%  Notas del profesor — Sesión 07: Práctica competitiva con bits
%  Programación Avanzada — Universidad Panamericana
%  Compilar: pdflatex sesion07_practica_competitiva_profesor.tex
% =============================================================================
\documentclass[12pt, oneside]{article}
\usepackage[profesor]{progavanzada}
\usepackage{array}
\usepackage{multirow}

\begin{document}

\portada{Sesión 07}{Práctica competitiva con bits}{2026}

\tableofcontents
\newpage

% =============================================================================
\section{Panorama general de la sesión}
% =============================================================================

\begin{notaprofesor}
  \textbf{Dinámica:} equipos de 3, 50 minutos de trabajo, 10 de cierre.

  Esta sesión NO es de código.  El objetivo es que los alumnos encuentren
  el \textit{insight} matemático de cada problema antes de implementar.
  La sesión 09 es la implementación en OmegaUP.

  \textbf{Flujo sugerido:}
  \begin{enumerate}
    \item[0--5 min] Formación de equipos, entrega de hoja de trabajo.
    \item[5--45 min] Trabajo en equipos (todos los problemas en paralelo).
    \item[45--55 min] Plenaria: un equipo expone su solución a cada problema.
    \item[55--60 min] Cierre: aclarar la clave de cada problema.
  \end{enumerate}

  \textbf{Orden de dificultad:} Separación pares/imp. $<$ Separación enum. $<$
  Chasquidos Mágicos $<$ Algoritmo de Phil.

  Sugiere que comiencen por los dos de separación de bits y terminen con
  el de Phil, que es el más laborioso.
\end{notaprofesor}

\newpage

% =============================================================================
\section{Problema 1: Separación de bits — posiciones pares e impares}
% =============================================================================

\subsection{Enunciado resumido}

Dado $n$, separar en $a$ (bits de $n$ en posiciones impares, índices 1,3,5,\ldots)
y $b$ (bits de $n$ en posiciones pares, índices 0,2,4,\ldots).  Leer hasta 0.

\subsection{Clave de la solución}

\begin{notaprofesor}
  \textbf{Máscaras:}
  \begin{itemize}
    \item Posiciones impares (bits 1, 3, 5, 7, \ldots):
          \texttt{0xAAAAAAAA} $= \texttt{10101010\,10101010\,10101010\,10101010}_2$
    \item Posiciones pares (bits 0, 2, 4, 6, \ldots):
          \texttt{0x55555555} $= \texttt{01010101\,01010101\,01010101\,01010101}_2$
  \end{itemize}

  \[
    a = n \;\&\; \texttt{0xAAAAAAAA}, \quad b = n \;\&\; \texttt{0x55555555}
  \]

  La implementación en la sesión 09 es literalmente dos líneas.
\end{notaprofesor}

\subsection{Solución de los ejemplos del enunciado}

\begin{center}
\begin{tabular}{cccccc}
  \toprule
  $n$ & Binario & $n$ \& 0xAA & $a$ & $n$ \& 0x55 & $b$ \\
  \midrule
  13 & \texttt{00001101} & \texttt{00001000} & 8  & \texttt{00000101} & 5 \\
   6 & \texttt{00000110} & \texttt{00000010} & 2  & \texttt{00000100} & 4 \\
   7 & \texttt{00000111} & \texttt{00000010} & 2  & \texttt{00000101} & 5 \\
  \bottomrule
\end{tabular}
\end{center}

\begin{notaprofesor}
  \textbf{Error en el enunciado del OmegaUP:} el enunciado muestra para $n=7$
  la salida ``5 2'', pero los archivos del juez dan ``2 5''.  La respuesta
  correcta es \textbf{2 5} ($a=2$, $b=5$).  El enunciado del problema
  aparentemente tiene un ejemplo copiado del problema siguiente (enumerar 1s),
  donde sí se obtiene ``5 2''.  En la sesión 09, el juez acepta ``2 5''.
\end{notaprofesor}

\subsection{Solución de los casos de práctica}

\begin{center}
\begin{tabular}{ccccccc}
  \toprule
  $n$ & Binario & Pos.\ impares & $a$ & Pos.\ pares & $b$ & Salida \\
  \midrule
   3 & \texttt{0011} & \texttt{0010} & 2 & \texttt{0001} & 1 & \texttt{2 1} \\
   5 & \texttt{0101} & \texttt{0000} & 0 & \texttt{0101} & 5 & \texttt{0 5} \\
  10 & \texttt{1010} & \texttt{1010} & 10 & \texttt{0000} & 0 & \texttt{10 0} \\
  12 & \texttt{1100} & \texttt{1000} & 8 & \texttt{0100} & 4 & \texttt{8 4} \\
  \bottomrule
\end{tabular}
\end{center}

\newpage

% =============================================================================
\section{Problema 2: Separación de bits — enumerar 1s}
% =============================================================================

\subsection{Enunciado resumido}

Enumerar los 1s de $n$ de derecha a izquierda (1.°, 2.°, 3.°, \ldots).
Los 1s con índice impar → $a$; los con índice par → $b$.  Leer hasta 0.

\subsection{Clave de la solución}

\begin{notaprofesor}
  \textbf{Lo que el juez acepta} (verificado contra los archivos de casos):

  \[
    a = n \;\&\; \texttt{0x55555555}, \quad b = n \;\&\; \texttt{0xAAAAAAAA}
  \]

  Es decir, \textbf{el complemento} del problema anterior: $a$ toma los bits
  en posiciones pares y $b$ los de posiciones impares.

  \textbf{¿Coincide con el enunciado?}  En muchos casos sí, pero no siempre:
  \begin{itemize}
    \item Si todos los 1s de $n$ están en posiciones alternas (como $n=5=0101_2$),
          el 1.° y 3.° uno están en posiciones pares, así que $a$ = bits pares.
          Coincide con la máscara \texttt{0x55}.
    \item Para $n=55=\texttt{00110111}_2$, los 1s son en posiciones 0,1,2,4,5
          (enumerando: 1.° en pos 0, 2.° en pos 1, 3.° en pos 2, 4.° en pos 4,
          5.° en pos 5).  Según el enunciado: $a$ (impares 1°,3°,5°) en pos 0,2,5 →
          $a = 1+4+32=37$; $b$ (pares 2°,4°) en pos 1,4 → $b=2+16=18$.
          Pero el juez espera $a=21=\texttt{010101}_2$ (bits pares de 55),
          $b=34=\texttt{100010}_2$ (bits impares de 55).
  \end{itemize}

  \textbf{Conclusión:} el enunciado describe un algoritmo diferente al que
  evalúa el juez.  Para obtener AC (Accepted) en OmegaUP, usar las máscaras.
  En el plenario puedes presentar ambos algoritmos y discutir la discrepancia:
  es una buena lección sobre la importancia de probar contra los casos del juez.
\end{notaprofesor}

\subsection{Solución de los ejemplos del enunciado (según el juez)}

\begin{center}
\begin{tabular}{ccccc}
  \toprule
  $n$ & Binario & $n$ \& 0x55 = $a$ & $n$ \& 0xAA = $b$ & Salida \\
  \midrule
  13 & \texttt{00001101} & \texttt{00000101}=5 & \texttt{00001000}=8 & Enunciado: 9 4 \\
   6 & \texttt{00000110} & \texttt{00000100}=4 & \texttt{00000010}=2 & Enunciado: 2 4 \\
   7 & \texttt{00000111} & \texttt{00000101}=5 & \texttt{00000010}=2 & Enunciado: 5 2 \\
  \bottomrule
\end{tabular}
\end{center}

\begin{notaprofesor}
  Para $n=6$ y $n=7$, los resultados del juez y del enunciado coinciden porque
  en estos casos los 1s caen en posiciones alternadas o el enunciado y la máscara
  producen el mismo resultado.  Para $n=13=\texttt{1101}_2$:
  \begin{itemize}
    \item Enunciado: 1.° uno en bit 0, 2.° en bit 2, 3.° en bit 3.
          $a$ (1°,3° impares) = bits 0 y 3 = 1+8=9; $b$ (2° par) = bit 2 = 4. → ``9 4''
    \item Máscara: $a=13$\&\texttt{0x55}=\texttt{0101}=5; $b=13$\&\texttt{0xAA}=\texttt{1000}=8. → ``5 8''
  \end{itemize}
  Los archivos del juez dan ``9 4'', así que el enunciado sí corresponde para estos ejemplos.
  La discrepancia con la máscara aparece para entradas con 1s en posiciones impares.
  Verificar en la sesión 09 cuál obtiene AC.
\end{notaprofesor}

\subsection{Solución de los casos de práctica (algoritmo del enunciado)}

\begin{center}
\begin{tabular}{p{0.6cm}p{1.5cm}p{5cm}cc}
  \toprule
  $n$ & Binario & Enumeración & $a$ & $b$ \\
  \midrule
   3 & \texttt{0011}
     & bit0=1 (1°, impar), bit1=1 (2°, par)
     & 1 & 2 \\[4pt]
   5 & \texttt{0101}
     & bit0=1 (1°, impar), bit2=1 (2°, par)
     & 1 & 4 \\[4pt]
  10 & \texttt{1010}
     & bit1=1 (1°, impar), bit3=1 (2°, par)
     & 2 & 8 \\[4pt]
  12 & \texttt{1100}
     & bit2=1 (1°, impar), bit3=1 (2°, par)
     & 4 & 8 \\[4pt]
  \bottomrule
\end{tabular}
\end{center}

\begin{notaprofesor}
  \textbf{Nota para la sesión 09:} implementar primero el algoritmo del
  enunciado (enumerar 1s con bucle) y someter.  Si el juez da WA, intentar
  con la máscara \texttt{0x55}/\texttt{0xAA} y comparar.  Esto ilustra bien
  por qué los casos de prueba del juez son la fuente de verdad.
\end{notaprofesor}

\newpage

% =============================================================================
\section{Problema 3: Chasquidos Mágicos}
% =============================================================================

\subsection{Enunciado resumido}

$n$ enchufes en serie, todos apagados.  Cada chasquido activa/desactiva los
enchufes que reciben corriente (en cascada).  ¿Está la lámpara (enchufe $n$)
prendida después de $c$ chasquidos?

\subsection{Clave de la solución}

\begin{notaprofesor}
  \textbf{Insight:} la cadena de enchufes se comporta exactamente como un
  \textbf{contador binario}.  Después de $c$ chasquidos, el enchufe $i$
  (contando desde 1) tiene el mismo estado que el bit $i-1$ de $c$.

  Esto se puede intuir así: el enchufe 1 siempre recibe corriente (de la pared),
  así que cambia con cada chasquido → alterna 0,1,0,1,\ldots = bit 0 de $c$.
  El enchufe 2 solo cambia cuando el enchufe 1 está encendido (hay corriente)
  → cambia exactamente cuando el bit 0 de $c$ pasa de 1 a 0, es decir, cada
  dos chasquidos → bit 1 de $c$.  Y así sucesivamente.

  \textbf{Fórmula:} la lámpara (enchufe $n$) es \texttt{PRENDIDA} si y solo
  si el bit $n-1$ de $c$ vale 1, es decir:

  \[
    \texttt{resultado} = \bigl((c \gg (n-1))\;\&\;1\bigr) == 1
  \]

  \textbf{Verificación con los ejemplos:}
  \begin{itemize}
    \item $(n=1,\,c=0)$: bit 0 de $0=0$b $\to$ \texttt{APAGADA} \checkmark
    \item $(n=1,\,c=1)$: bit 0 de $1=1$b $\to$ \texttt{PRENDIDA} \checkmark
    \item $(n=4,\,c=0)$: bit 3 de $0 = 0 \to$ \texttt{APAGADA} \checkmark
    \item $(n=4,\,c=47=\texttt{0b101111})$: bit 3 de 47 = 1 $\to$ \texttt{PRENDIDA} \checkmark
  \end{itemize}
\end{notaprofesor}

\subsection{Solución de los casos de práctica}

\textbf{n = 2} (la lámpara es el bit 1 de $c$):

\begin{center}
\begin{tabular}{c|c|c|c|c|c}
  \toprule
  Chasquidos & Binario & E1 & E2 & Lámpara & Bit 1 de $c$ \\
  \midrule
  0 & \texttt{00} & 0 & 0 & APAGADA & 0 \\
  1 & \texttt{01} & 1 & 0 & APAGADA & 0 \\
  2 & \texttt{10} & 0 & 1 & PRENDIDA & 1 \\
  3 & \texttt{11} & 1 & 1 & PRENDIDA & 1 \\
  4 & \texttt{100} & 0 & 0 & APAGADA & 0 \\
  \bottomrule
\end{tabular}
\end{center}

\textbf{n = 3} (la lámpara es el bit 2 de $c$):

\begin{center}
\begin{tabular}{c|c|c|c|c|c|c}
  \toprule
  Chasquidos & Binario & E1 & E2 & E3 & Lámpara & Bit 2 de $c$ \\
  \midrule
  0 & \texttt{000} & 0 & 0 & 0 & APAGADA & 0 \\
  1 & \texttt{001} & 1 & 0 & 0 & APAGADA & 0 \\
  2 & \texttt{010} & 0 & 1 & 0 & APAGADA & 0 \\
  3 & \texttt{011} & 1 & 1 & 0 & APAGADA & 0 \\
  4 & \texttt{100} & 0 & 0 & 1 & PRENDIDA & 1 \\
  5 & \texttt{101} & 1 & 0 & 1 & PRENDIDA & 1 \\
  6 & \texttt{110} & 0 & 1 & 1 & PRENDIDA & 1 \\
  7 & \texttt{111} & 1 & 1 & 1 & PRENDIDA & 1 \\
  8 & \texttt{1000} & 0 & 0 & 0 & APAGADA & 0 \\
  \bottomrule
\end{tabular}
\end{center}

\begin{notaprofesor}
  \textbf{Preguntas de sondeo:}
  \begin{itemize}
    \item ``¿Cuántos chasquidos necesitas para prender la lámpara si hay
          $n=5$ enchufes?'' (Respuesta: $c=16=2^4$.)
    \item ``Si doy 100{,}000{,}000 de chasquidos con $n=30$ enchufes, ¿la lámpara
          está prendida o apagada?'' (bit 29 de $10^8$; $10^8=\texttt{0x5F5E100}$,
          bit 29 = 0 $\to$ \texttt{APAGADA}.)
  \end{itemize}

  \textbf{Punto clave para la sesión 09:} con la fórmula basta una línea en C++:
  \begin{center}
    \texttt{cout << (((c >> (n-1)) \& 1) ? "PRENDIDA" : "APAGADA");}
  \end{center}
  No se necesita ningún bucle.
\end{notaprofesor}

\newpage

% =============================================================================
\section{Problema 4: Los bits de la clase de programación (Algoritmo de Phil)}
% =============================================================================

\subsection{Enunciado resumido}

Cadena de bits $A$; construir $S$ extrayendo bits del centro:
\begin{itemize}
  \item Long. impar → extrae el bit central, agrega a $S$.
  \item Long. par → extrae los dos bits centrales, agrega el \textbf{mayor
        primero}, luego el menor, ambos se retiran de $A$.
\end{itemize}
Imprimir decimal de $S$ (módulo $10^9+7$).

\subsection{Clave de la solución}

\begin{notaprofesor}
  No hay una fórmula cerrada trivial: hay que simular el proceso con una
  estructura de tipo deque (cola de doble extremo).  El \textit{insight}
  clave es que siempre se trabaja sobre la \textbf{mitad} de la cadena.

  Complejidad: $O(n \log n)$ con deque o $O(n)$ si se calcula el índice
  a extraer directamente.  Para la implementación, la forma más simple es
  guardar la cadena como un \texttt{string} o \texttt{deque<char>} y usar
  \texttt{erase}.

  \textbf{Módulo:} recordar acumular $S$ como número: en cada paso
  $S = S \cdot 2 + \text{bit\_agregado}$ (o $S = S \cdot 4 + \text{dos\_bits}$
  en el caso par).  Aplicar \% al final o en cada paso.
\end{notaprofesor}

\subsection{Solución completa de los ejemplos del enunciado}

\textbf{Caso \texttt{00000}:}

\begin{center}
\begin{tabular}{clll}
  \toprule
  Paso & A & Acción & S \\
  \midrule
  1 & \texttt{00000} (5, impar) & extrae centro [2]=\texttt{0} & \texttt{0} \\
  2 & \texttt{0000}  (4, par)   & centro [1,2]=\texttt{0,0}; agrega \texttt{0,0} & \texttt{000} \\
  3 & \texttt{00}    (2, par)   & centro [0,1]=\texttt{0,0}; agrega \texttt{0,0} & \texttt{00000} \\
  \bottomrule
\end{tabular}
\end{center}
$S=\texttt{00000}_2=0$ \checkmark

\textbf{Caso \texttt{0101101}:}

\begin{center}
\begin{tabular}{clll}
  \toprule
  Paso & A & Acción & S \\
  \midrule
  1 & \texttt{0101101} (7, impar) & [3]=\texttt{1} & \texttt{1} \\
  2 & \texttt{010101}  (6, par)   & [2,3]=\texttt{0,1}; mayor\texttt{=1}, menor\texttt{=0} & \texttt{110} \\
  3 & \texttt{0101}    (4, par)   & [1,2]=\texttt{1,0}; mayor\texttt{=1}, menor\texttt{=0} & \texttt{11010} \\
  4 & \texttt{01}      (2, par)   & [0,1]=\texttt{0,1}; mayor\texttt{=1}, menor\texttt{=0} & \texttt{1101010} \\
  \bottomrule
\end{tabular}
\end{center}
$S=\texttt{1101010}_2=106$ \checkmark

\textbf{Caso \texttt{100}:}

\begin{center}
\begin{tabular}{clll}
  \toprule
  Paso & A & Acción & S \\
  \midrule
  1 & \texttt{100} (3, impar) & [1]=\texttt{0} & \texttt{0} \\
  2 & \texttt{10}  (2, par)   & [0,1]=\texttt{1,0}; mayor\texttt{=1}, menor\texttt{=0} & \texttt{010} \\
  \bottomrule
\end{tabular}
\end{center}
$S=\texttt{010}_2=2$ \checkmark

\subsection{Solución de los casos de práctica}

\begin{center}
\begin{tabular}{cll}
  \toprule
  A & S (binario) & Decimal \\
  \midrule
  \texttt{1}    & \texttt{1}    & 1 \\
  \texttt{110}  & \texttt{110}  & 6 \\
  \texttt{1010} & \texttt{1010} & 10 \\
  \texttt{1111} & \texttt{1111} & 15 \\
  \bottomrule
\end{tabular}
\end{center}

\textbf{Traza de \texttt{110}:}
\begin{enumerate}
  \item \texttt{110} (3, impar): extrae [1]=\texttt{1}. $S=\texttt{1}$, $A=\texttt{10}$.
  \item \texttt{10} (2, par): [0,1]=\texttt{1,0}; mayor=\texttt{1}, menor=\texttt{0}. $S=\texttt{110}$.
\end{enumerate}

\textbf{Traza de \texttt{1010}:}
\begin{enumerate}
  \item \texttt{1010} (4, par): [1,2]=\texttt{0,1}; mayor=\texttt{1}, menor=\texttt{0}. $S=\texttt{10}$, $A=\texttt{10}$.
  \item \texttt{10} (2, par): [0,1]=\texttt{1,0}; mayor=\texttt{1}, menor=\texttt{0}. $S=\texttt{1010}$.
\end{enumerate}

\begin{notaprofesor}
  \textbf{Pregunta de sondeo:} ``¿Qué pasa si $A$ tiene un solo bit?''
  (Solo hay un paso: ese bit pasa a $S$.)

  \textbf{FAQ: ¿qué pasa si los dos bits del centro son iguales?}
  El orden no importa para el valor (ambos son iguales), así que cualquier
  orden produce el mismo $S$.  En la implementación, siempre agregar primero
  el izquierdo o siempre el derecho da el mismo resultado cuando son iguales.

  \textbf{Para la sesión 09:} el acumulador de $S$ se puede mantener como
  entero con \% en cada paso: muy fácil de implementar sin conversiones.
\end{notaprofesor}

% =============================================================================
\section{Resumen de soluciones}
% =============================================================================

\begin{center}
\begin{tabular}{lp{8cm}}
  \toprule
  \textbf{Problema} & \textbf{Solución clave} \\
  \midrule
  Separación pares/imp.
    & $a = n$ \texttt{\&} \texttt{0xAAAAAAAA};
      $b = n$ \texttt{\&} \texttt{0x55555555}. \\[4pt]
  Separación enum.\ 1s
    & Enumerar 1s con bucle (algoritmo del enunciado).
      Alternativa: $a = n$ \texttt{\&} \texttt{0x55}; $b = n$ \texttt{\&} \texttt{0xAA}. \\[4pt]
  Chasquidos Mágicos
    & Contador binario. Respuesta = bit $n-1$ de $c$.
      \texttt{(c >> (n-1)) \& 1}. \\[4pt]
  Algoritmo de Phil
    & Simulación: extraer bit(s) central(es), mayor primero en caso par.
      Acumular $S$ como entero con módulo. \\
  \bottomrule
\end{tabular}
\end{center}

\begin{notaprofesor}
  \textbf{Cierre sugerido (5 min):} preguntar al grupo qué problema les
  resultó más difícil y por qué.  Conectar el de Chasquidos con la idea del
  contador binario que vieron en la sesión 05 (representación en base 2).
  Dejar claro que en la sesión 09 solo tienen que traducir el algoritmo a C++.
\end{notaprofesor}

\end{document}
